% Modernised reading — fonds Alexandre Grothendieck, shelfmark 100
% (pages 1-52, the whole folder).
%
% This is an INTERPRETATION, not a transcription. It re-reads the pages in
% today's notation and vocabulary, states the hypotheses the manuscript leaves
% implicit, and drops the critical apparatus so that the mathematics reads
% straight through. Everything the brackets and underlines carried moves into a
% footnote. It is held to being mathematically correct as it stands; where the
% manuscript is loose or mistaken, the true statement is what appears here and
% the footnote says what the page has. In any dispute of substance the
% transcription governs: whoever cites Grothendieck must cite batch-01.fr.tex
% (pp. 1-20), batch-02.fr.tex (pp. 21-40) or batch-03.fr.tex (pp. 41-52).
%
% Pass: Opus 5.5 (claude-opus-5-5), 2026-10-10 — first pass, unchecked against the pages by a human.
% The folder was read whole, all three batches; every batch is transcribed.
%
% What the folder is. Notes in blue ink (pp. 2-44, 51-52) and a four-leaf
% typescript (pp. 47-50) around one statement: over a locally noetherian S,
% with U open and every residue characteristic on T = S - U equal to p, an
% abelian scheme over U extends to S as soon as its p-divisible group does,
% and the extension is unique — the functor A |-> (A_U, A[p^oo], id) is fully
% faithful when U meets every connected component, an equivalence under
% geometric unibranchness on T and finiteness of the normalisation (p. 2).
% The other leaves of the folder — an English typescript on quotients by
% group actions (pp. 5-15, 21-37), EGA IV errata (pp. 11, 17, 41-45), the
% cover of Publ. Math. IHES 32 (p. 19), two letters (pp. 29, 39) — are the
% paper he wrote on or kept with it; only his pencil on the first is noted.
%
% Notation fixed for the whole document. A[p^oo] for what he writes T_p(A)
% (the p-divisible group, not the Tate module); triples (A_U, M, phi) with
% phi : M|_U -> A_U[p^oo] an isomorphism; S_u for the "lieu d'algébricité"
% of a homomorphism u of p-divisible groups (his T on pp. 3 and 20, renamed
% because T is already the complement of U); \mathcal{M} for the formal
% moduli of p. 28 (his M, renamed because M is the p-divisible group); in
% the typescript, H = (H_n) a p-divisible group over the trait S, G_eta[p^oo]
% for his wavy-underlined G_eta. S', S'' = S' x_S S' are reused by the pages
% for different auxiliary bases in different sections; each section recalls
% which.
%
% Substantive departures from the pages, each footnoted where it occurs:
% — p. 2: « U rencontre toutes les composantes connexes (i.e. U dense, i.e.
%   prof_T(S) >= 1) » joins three inequivalent conditions; the reading keeps
%   the first, which is what pp. 3 and 20 actually use and suffices.
% — p. 3: the Complément holds at points of residue characteristic l (= p),
%   by Serre-Tate; the reading says so; the existence of the closed
%   subscheme rests on « Invariants 1.2 », not identified.
% — p. 4: « M gpe p-div. sur T » read « sur S »; X read S; the mark on T_p
%   (\natural in the transcription) dropped.
% — p. 8: « prouver que c'est un isom. » — for S geometrically unibranch and
%   not normal the argument gives at most a finite, birational, radicial
%   S' -> S; the reading says so.
% — p. 22: « S'_1 -> S' fini radiciel » read S'_1 -> S.
% — p. 28: W(k) is the right base only for k perfect (Cohen ring otherwise).
% — p. 30: « quotient intègre B de A » read « de R ».
% — p. 34: « S'_i -> S immersion fermée » needs S' -> S unramified, true for
%   the S' built on p. 32 (finite unions of the net P^alpha), not for an
%   arbitrary finite S'; the first Lemma's gluing is footnoted (pinching).
% — p. 34/36: the S_red Lemma has « lisse » struck; the proof of p. 36 needs
%   it; the reading states the Lemma for smooth X^ (enough for abelian ones).
% — p. 36: the Proposition is false for a non-complete local ring
%   (countably many height-one primes); « complet », struck on the page, is
%   restored, as the proof needs it.
% — p. 47: Lemma 1's v(m,g) = u(m) + g is a homomorphism only if G is
%   commutative (his pencil « G commutatif ? »); supplied.
% — p. 51: phi factors through the centre of Z, not through Centr(Z) as
%   written; case 2 needs the boxed hypothesis and G_red.
% — p. 52: « H -i-> K » read H -> G; « modèle de Néron » of H_K is of
%   finite type because G_K has none of G_m either (footnoted).
%
% Modern names given and footnoted as ours: Barsotti-Tate group (his own
% term from 1970, absent from the pages), Serre-Tate theorem, Tate's theorem
% and de Jong's extension to characteristic p (1998), p-adic criterion of
% good reduction, Grothendieck existence theorem, Néron mapping property,
% smoothening, Weil restriction, semistable reduction theorem (SGA 7 IX),
% Chevalley's structure theorem, Baire category, Cohen structure theorem,
% Ferrand pinching, Anantharaman's quotients.
%
% What the folder announces and never establishes: the essential
% surjectivity of p. 2 in general (b) closes on a struck sketch, e) rests on
% a Lemma he marks « Douteux », p. 16 breaks off, p. 18 stops on « Alors »);
% the proof of the Lemma of p. 3 (a reference); Corollary 3 of p. 24 (Hom of
% p-divisible groups formally unramified); most of the proof of p. 36;
% p. 40's equivalences (« procédant comme plus haut »); the descent remarks
% of pp. 42-44 (fragments); Lemma 3 of p. 49 (cited from « le séminaire de
% l'an dernier »); Tate's theorem when the generic point has characteristic
% p (the typescript says so itself); the typescript's Corollary (no proof).
\documentclass[11pt,a4paper]{article}
\input{../preamble/grothendieck}

\folder{100}
\pages{1}{52}
\dating{[à partir de 1967-1968]}
\watermark{Édition de démonstration\\interprétation personnelle de l'œuvre}
\foldertitle{Bonne réduction des variétés abéliennes via bonne réduction de Barsotti-Tate Tp (A) (base quelconque) : notes manuscrites (s.d.), copies de tapuscrit annoté (s.d.) — lecture modernisée du dossier entier}

\begin{document}

\section*{Résumé}

\begin{resume}
Une \emph{variété abélienne} est la généralisation en dimension quelconque
d'une courbe elliptique : un objet géométrique qui est en même temps un
groupe commutatif. Quand ses équations sont à coefficients entiers, on peut
les réduire modulo un nombre premier ; si l'objet réduit est encore une
variété abélienne, on dit qu'il y a \emph{bonne réduction}. Plus
généralement, on se donne une famille de variétés abéliennes paramétrée par
un ouvert $U$ d'un espace $S$, et l'on demande si elle se prolonge sans
dégénérer au-dessus des points de $S$ qui manquent.

Pour le savoir, on regarde un échantillon de la famille : ses points de
torsion. Les points d'ordre une puissance de $p$ forment une suite de petits
groupes finis emboîtés, qu'on appelle aujourd'hui un \emph{groupe
$p$-divisible} ou \emph{groupe de Barsotti--Tate}. Lorsque $p$ est inversible,
ces petits groupes sont de simples ensembles de points sur lesquels agit un
groupe de Galois, et le critère classique de Néron, Ogg et Šafarevič dit que
la bonne réduction se lit sur eux. Lorsque $p$ est justement la
caractéristique des points qui manquent, ce sont des objets beaucoup plus
fins, qui portent une partie infinitésimale ; et l'idée de ces feuillets est
qu'ils en savent assez pour décider de tout : \emph{la famille se prolonge si
et seulement si son groupe $p$-divisible se prolonge}, et le prolongement est
unique. Un thermomètre placé à l'intérieur d'un objet ne donne d'habitude
qu'une information partielle ; ici, à condition de regarder la torsion de la
bonne caractéristique, il donne tout.

Le titre du dossier dit l'ambition : une \emph{base quelconque}. Sur une base
de dimension un --- un trait, l'analogue d'un disque autour d'un point ---,
l'énoncé découle d'un théorème de Tate, et un tapuscrit annoté rangé dans le
dossier en donne la démonstration. Grothendieck veut passer de là à une base
localement noethérienne arbitraire. Les instruments sont ceux de sa géométrie
algébrique : la \emph{descente}, qui recolle un objet construit sur un
revêtement de la base ; la \emph{géométrie formelle}, qui construit d'abord
l'objet sur des voisinages infinitésimaux de plus en plus gros, puis le rend
algébrique ; les \emph{espaces de modules}, qui transforment une famille de
variétés abéliennes en une application vers un espace qui les classe ; et,
de façon inattendue, un argument de Baire : une réunion dénombrable de
fermés rares ne remplit pas un anneau local complet.

Le dossier n'arrive pas au bout. Le théorème de pleine fidélité (l'unicité du
prolongement) est démontré ; l'existence est ramenée, par une suite de
réductions dont certaines ne sont qu'esquissées, à des énoncés de descente et
d'algébrisation qui occupent la seconde moitié des notes et qui sont souvent
intéressants pour eux-mêmes. Le tapuscrit avoue la faiblesse de l'édifice :
le théorème de Tate sur lequel tout repose n'était démontré qu'en
caractéristique nulle. Les feuillets s'achèvent sur un lemme de sa main qui
supprime une hypothèse du tapuscrit, et sur « On gagne ! ».

\keywords{abelian scheme, Barsotti-Tate group, good reduction, Serre-Tate theorem, Tate's theorem on p-divisible groups, fpqc descent, effective epimorphism, formal moduli, algebraization of formal schemes, Grothendieck existence theorem, Néron model, semistable reduction, Weil restriction, smoothening, Baire category theorem}
\end{resume}

\section*{Le fil du dossier, et les conventions}

\subsection*{Les stations}

L'inventaire date le dossier « à partir de 1967-1968 » et le range dans le
groupe « Géométrie et topologie combinatoire » (dossiers 69 à 102, daté
1976-[vers 1986]), dont le sujet n'est pas le sien.\footnote{Les lectures
modernisées des dossiers 81, 88 et 89 du même groupe portent sur les cartes,
les dessins d'enfants et des structures combinatoires finies ; aucune ne
rencontre celui-ci, et on ne les cite plus. Le dossier contient aussi, au
verso ou entre les feuillets de sa main, des pages qui ne sont pas de lui : un
tapuscrit anglais sur les quotients d'un schéma par un groupe (pages 5 à 15,
21 à 37), des listes d'errata d'EGA IV (pages 11, 17, 41 à 45), la couverture
du numéro 32 des \emph{Publications mathématiques} (page 19), une lettre de
L.~Illusie datée du 24 février 1968 (page 29) et une lettre à Dieudonné datée
du 7 avril 1967 (page 39). On ne les résume pas ; on ne tire rien de leurs
dates.}

\begin{enumerate}
  \item \emph{L'énoncé} (page 2, repris page 4) : le foncteur
  $A \mapsto (A_U, A[p^\infty], \mathrm{id})$ est pleinement fidèle, et une
  équivalence sous des hypothèses d'unibranchement.
  \item \emph{La pleine fidélité} (pages 3 et 4, reprise page 20) : le lieu où
  un homomorphisme de groupes $p$-divisibles est algébrique est un fermé qui
  contient un voisinage de chacun de ses points de caractéristique $p$.
  \item \emph{L'essentielle surjectivité}, par réductions successives (pages 6
  à 18) : $S$ réduit, $S$ géométriquement unibranche, $S$ normal, $S$ réduit
  par descente le long du normalisé, cas général par recollement.
  \item \emph{La descente par un épimorphisme propre} (pages 21 à 24) : pleine
  fidélité de la descente pour les schémas abéliens.
  \item \emph{L'effectivité, et l'algébrisation des schémas abéliens formels}
  (pages 26 à 38) : réduction à une base locale complète, construction formelle
  par les modules formels, un théorème d'algébrisation à trois conditions
  équivalentes, et l'argument de Baire qui le termine.
  \item \emph{Le cas d'un schéma propre et plat quelconque} (pages 40 à 44) :
  une remarque, et des remarques de descente en fragments.
  \item \emph{Le cas d'un trait} (pages 47 à 52) : le tapuscrit « Modèles de
  Néron et groupes $p$-divisibles », puis un lemme de sa main qui en supprime
  une hypothèse.
\end{enumerate}

Les stations 4 à 6 sont les outils que réclame la station 3 : la réduction b)
de la page 8 invoque le cas d'un trait, qui est la station 7 ; les réductions
c) et d) des pages 10 à 14 demandent que les morphismes finis surjectifs
soient de descente effective pour les schémas abéliens, ce que les pages 22 à
38 entreprennent d'établir ; la réduction a) de la page 6 a besoin du lemme
d'algébrisation « modulo le nilradical » des pages 34 à 36. Les feuillets ne
font eux-mêmes aucun de ces renvois, sinon par les mots « Raynaud » et
« Tate » qui reviennent d'un bout à l'autre ; ces liens sont notre lecture.

\subsection*{Conventions, valables pour tout le dossier}

$S$ est un schéma localement noethérien, $U \subset S$ un ouvert, $T = S - U$ le
fermé complémentaire, et $p$ un nombre premier tel que tout point $t \in T$ ait
un corps résiduel $k(t)$ de caractéristique $p$. Pour un schéma abélien $A$, on
note $A[p^\infty] = (A[p^n])_n$ son groupe $p$-divisible.\footnote{Il écrit
$T_p(A)$, la notation du module de Tate, pour le groupe $p$-divisible tout
entier ; sur une base où $p$ n'est pas inversible, le module de Tate ne suffit
pas, et c'est bien du groupe $p$-divisible qu'il s'agit (les « triples » de la
page 2 contiennent un « groupe $p$-divisible $M$ »). Le titre de l'inventaire
dit « Barsotti-Tate », nom que Grothendieck lui-même donnera à ces groupes
(exposé au congrès de Nice, 1970) ; les feuillets disent « groupe
$p$-divisible », le terme de Tate.} Un homomorphisme
$v : A[p^\infty] \to B[p^\infty]$ est dit \emph{algébrique} s'il est induit par
un homomorphisme de schémas abéliens $u : A \to B$ ; on note $S_v$ le lieu
où il l'est (page 3).\footnote{Il appelle ce lieu $T$, comme le fermé
complémentaire de $U$ ; on le renomme.} $\hat{\ }$ désigne une complétion
formelle, $S_{\mathrm{red}}$ le sous-schéma réduit, $\mathcal{O}_{S,s}$ l'anneau
local en $s$ et $S(s) = \mathrm{Spec}\,\mathcal{O}_{S,s}$. Les lettres $S'$,
$S''= S' \times_S S'$, $S''' = S' \times_S S' \times_S S'$ désignent, d'une
section à l'autre, des bases auxiliaires différentes, comme sur les pages ; on
rappelle chaque fois laquelle.

\pagerange{2}{2}
\section*{L'énoncé (page 2)}

\textbf{Théorème.} Soient $C$ la catégorie des schémas abéliens sur $S$, et
$C'$ celle des triples $(A_U, M, \varphi)$, où $A_U$ est un schéma abélien sur
$U$, $M$ un groupe $p$-divisible sur $S$ et
$\varphi : M|_U \xrightarrow{\ \sim\ } A_U[p^\infty]$ un isomorphisme. Soit
\[
  F : C \longrightarrow C', \qquad
  F(A) = \bigl(A \times_S U,\ A[p^\infty],\ \mathrm{id}\bigr).
\]
\begin{enumerate}
  \item[1)] Si $U$ rencontre toutes les composantes connexes de $S$, $F$ est
  pleinement fidèle.\footnote{La page ajoute en interligne « (i.e.\ $U$ dense
  dans $S$, i.e.\ $\mathrm{prof}_T(S) \geq 1$) ». Les trois conditions ne sont
  pas équivalentes : la troisième (aucun point associé de $\mathcal{O}_S$ dans
  $T$, c'est-à-dire $U$ schématiquement dense) entraîne la deuxième, qui
  entraîne la première, et aucune réciproque n'est vraie. C'est la
  première --- la plus faible --- qu'utilise la démonstration de la page 3, et
  elle suffit ; la page 4 dit d'ailleurs que le cas schématiquement dense est
  le cas facile. Le dernier mot, « connexes », est mal formé.}
  \item[2)] Si de plus $S$ est géométriquement unibranche en tout point de
  $T$, et si le normalisé de $S_{\mathrm{red}}$ est fini sur $S$, $F$ est une
  équivalence de catégories.\footnote{Le mot qui suit « Supposons » est
  recouvert ; l'hypothèse est lue sur la ligne suivante (« $S'$ géom.\
  unibranche ») et sur l'en-tête de la page 6, où elle est écrite en
  entier : « Si $S$ géom.\ unibranche en les points de $S - U$ et si le
  normalisé de $S_{\mathrm{red}}$ est fini sur $S$ ». Isolé sous le cadre,
  « $\mathrm{prof}_T(S) \geq 2$ », sans attache.}
\end{enumerate}

Une condition de rechange, plus combinatoire, est essayée puis abandonnée.
D'abord, biffée : pour tout $t \in T$, deux composantes irréductibles du
localisé strict de $S$ en $t$ se rencontrent ailleurs qu'en son point fermé.
Puis, non biffée mais inachevée : pour un morphisme fini surjectif
$S' \to S$, on forme l'ensemble simplicial tronqué des composantes connexes
\[
  L_0 = \pi_0(S') \leftleftarrows L_1 = \pi_0(S'')
  \;\substack{\leftarrow\\[-2pt]\leftarrow\\[-2pt]\leftarrow}\;
  L_2 = \pi_0(S'''),
\]
et le sous-ensemble simplicial $K_\bullet$ des composantes dont l'image dans
$S$ rencontre $U$ ; la condition b$'$) demande que la plus petite partie de
$L_1$ contenant $K_1$ et stable par l'opération déduite des faces de $L_2$ soit
$L_1$ tout entier.\footnote{« stable par “troisième face” », lecture
incertaine. Le $S'$ renvoie aux « notations de a) », un alinéa a) qui est
biffé ; on comprend un revêtement fini de $S$ par un schéma géométriquement
unibranche, comme son normalisé. Le sens de la condition, qui est le nôtre :
une donnée de descente le long de $S' \to S$ est une donnée sur $S''$ ; si
elle n'est connue que sur l'image inverse de $U$, b$'$) dit que la
compatibilité sur $S'''$ suffit à la propager à toutes les composantes de
$S''$. La page s'arrête sur « Alors $F$ est une équivalence de
catégories\ldots ».}
C'est la forme que prendra la réduction d) des pages 12 et 14, où il faudra
que $U''$ soit dense dans $S''$.

\pagerange{3}{5}
\section*{La pleine fidélité (pages 3 et 4)}

La fidélité est connue : $u \mapsto u[p^\infty]$ est injectif, parce qu'un
homomorphisme de schémas abéliens nul sur les fibres est nul (rigidité), et
que sur un corps un homomorphisme qui tue tous les $A[p^n]$ se factorise par
toutes les multiplications par $p^n$, donc est nul, $\mathrm{Hom}(A,B)$ étant
un groupe abélien de type fini. Pour la plénitude, il faut montrer qu'un
homomorphisme $v : A[p^\infty] \to B[p^\infty]$ défini sur $S$ et algébrique
sur $U$ est algébrique sur $S$ ; l'unicité de l'homomorphisme $u$ qui
l'induit, par la fidélité, fera le reste.

\textbf{Lemme} (page 3). Soient $A$, $B$ deux schémas abéliens sur $S$ et
$v : A[p^\infty] \to B[p^\infty]$. Il existe un sous-schéma fermé $S_v$ de $S$
tel que, pour tout $S$-schéma $S'$, $v_{S'}$ soit induit par un homomorphisme
$A_{S'} \to B_{S'}$ si et seulement si $S' \to S$ se factorise par
$S_v$.\footnote{La page énonce le lemme pour un $u_\ell : T_\ell(A) \to
T_\ell(B)$ et un nombre premier $\ell$ ; ici $\ell = p$, et la page 20 reprend
l'énoncé avec $p$. Elle ne le démontre pas : « Cf Invariants 1.2
(démonstration analogue) ». Nous n'identifions pas ce texte. La page 20
ajoute que la formation de $S_v$ commute à tout changement de base
(localement noethérien), et met « (cf.\ Serre) » au-dessus de « plus
grand », lecture incertaine.}

\textbf{Complément.} Si $t \in S_v$ et si $k(t)$ est de caractéristique $p$,
alors $S_v$ contient un voisinage ouvert de $t$.

C'est le théorème de Serre--Tate qui le donne, bien que la page 3 ne le dise
pas.\footnote{La page 3 écrit « $T$ majore un voisinage ouvert de $t$ »,
et la page 20 répète l'énoncé. La justification est la nôtre ; le théorème de
Serre--Tate est nommé à la page 4, dans un argument biffé.} Sur
$S_n = \mathrm{Spec}\,\mathcal{O}_{S,t}/\mathfrak{m}_t^{n+1}$, où $p$ est
nilpotent, un homomorphisme de schémas abéliens est la même chose qu'un
homomorphisme de leurs réductions au point fermé et de leurs groupes
$p$-divisibles qui se correspondent ; comme $v_t$ est algébrique, $v_{S_n}$
l'est pour tout $n$. L'idéal de $S_v$ est donc nul dans le complété de
$\mathcal{O}_{S,t}$, donc dans $\mathcal{O}_{S,t}$, donc au voisinage de $t$.

\textbf{Corollaire} (page 3). Si toute composante connexe de $S$ rencontre $U$
et si $v|_U$ est algébrique, $v$ est algébrique.

\emph{Démonstration.} Le sous-schéma fermé $S_v$ contient $U$ (hypothèse), et
contient un voisinage de chacun de ses points situés dans $T$ (complément,
puisque ces points sont de caractéristique $p$). Son idéal est donc nul au
voisinage de chacun de ses points : $S_v$ est un sous-schéma à la fois ouvert
et fermé, une réunion de composantes connexes de $S$. Toutes rencontrent
$U \subset S_v$ ; donc $S_v = S$.\footnote{Avant cette conclusion, un essai
par les composantes irréductibles, encadré et biffé de deux longs traits. La
page souligne ou surligne tour à tour les $T$ de la conclusion, pour
distinguer, semble-t-il, l'ensemble sous-jacent et le sous-schéma.}

Cela démontre 1). La page 4 récrit l'énoncé --- triples, $T$ de
caractéristiques résiduelles toutes égales à $p$, \emph{spécialisation}
$U$ schématiquement dense ---,\footnote{Elle écrit « $M$ gpe $p$-div.\ sur
$T$ », qu'on lit « sur $S$ » comme à la page 2, et
$\varphi : T_p(A_U) \simeq M$, l'isomorphisme dans l'autre sens. Elle munit
$T_p$ d'un petit signe en exposant, à chaque occurrence, peut-être une
marque de complétion ; on ne le reprend pas. Elle note $X$ ce qui est
clairement $S$.} puis la démonstration : si $U$ est schématiquement dense,
cela résulte de la référence laissée en blanc « [\;] 1.2 » ; dans le cas
général, on remplace $S$ par l'adhérence schématique $S_0$ de $U$, sur
laquelle on dispose d'un homomorphisme $A_{S_0} \to B_{S_0}$ compatible avec
$v$, et l'on remonte de $S_0$ à $S$ par Serre--Tate le long des voisinages
infinitésimaux de $T$, puis par complétion et descente fidèlement
plate.\footnote{Cette fin d'argument est encadrée et biffée de deux traits ;
on n'en donne que le mouvement. Le corollaire de la page 3 la rend inutile.}

\emph{La marge de la page 4.} Un bloc cerné, lourdement repris, dit que pour la
pleine fidélité il suffit que $U$ soit dense, et évoque un contre-exemple
obtenu en recollant les spectres de deux anneaux de valuation discrète
complets, $U$ étant l'un des deux points maximaux ; il se termine sur « Y
a-t-il des contre-exemples avec $S$ de dim.\ 2 ? ».\footnote{L'objet du
contre-exemple ne se lit pas. Ce n'est pas la pleine fidélité : le corollaire
de la page 3 s'applique à ce $S$, qui est connexe. Sur deux traits recollés en
leurs points fermés, la base n'est pas géométriquement unibranche au point de
recollement, et l'essentielle surjectivité y est en défaut --- deux schémas
abéliens sur les deux traits dont les fibres spéciales ont des groupes
$p$-divisibles isomorphes sans être isomorphes ne se recollent pas, alors que
les triples se recollent. Cette lecture est la nôtre.}
Au verso (page 5), un feuillet d'un tapuscrit anglais sur les quotients par
un groupe porte quelques notes de lui au crayon.\footnote{Ce tapuscrit
n'est pas de lui ; on ne le résume pas. Ses notes au crayon (pages 5, 7, 9,
15 et 21) portent sur ce texte et non sur le dossier : une définition du
quotient $X/R$ par le faisceau des invariants et les ouverts invariants,
« I think sufficient if $\varphi$ is qu.\ cpct qu.\ sep. », une correction de
lettre, « Dire : \emph{propre} et libre ».}

\section*{L'essentielle surjectivité (pages 6 à 18)}

On part d'un triple $(A_U, M, \varphi)$ et l'on cherche $A$. L'hypothèse
est celle de 2) : $S$ géométriquement unibranche aux points de $T$, normalisé
de $S_{\mathrm{red}}$ fini sur $S$. La page la dédouble aussitôt en une
hypothèse de descente : pour $S' \to S$ convenable, l'image inverse $U'$
(resp.\ $U''$) de $U$ rencontre toutes les composantes connexes de $S'$
(resp.\ $S''$) --- de sorte que la pleine fidélité vaille sur $S'$ et $S''$.
On peut supposer $S = \mathrm{Spec}\,\Lambda$ affine.

\pagerange{6}{7}
\subsection*{a) Réduction au cas $S$ réduit (page 6)}

Soit $S_0 = S_{\mathrm{red}}$, et supposons trouvé sur $S_0$ un schéma abélien
$A_0$ qui résout le problème pour le triple induit. Soit $T_n$ le $n$-ième
voisinage infinitésimal de $T$ dans $S$ ; $p$ y est nilpotent, puisque tous
les points de $T$ sont de caractéristique $p$. Par Serre--Tate, $A_0$ et $M$
définissent un schéma abélien sur chaque $T_n$, donc un schéma abélien formel
sur le complété $\mathrm{Spf}\,\hat\Lambda$ de $S$ le long de $T$. Sa
restriction au-dessus de $S'_0 = (\mathrm{Spec}\,\hat\Lambda)_{\mathrm{red}}$
est algébrique --- c'est l'image inverse de $A_0$ ---, donc il est algébrique
sur $S' = \mathrm{Spec}\,\hat\Lambda$, puisque $S'_0 \to S'$ est une immersion
nilpotente : on obtient un schéma abélien $A'$ sur $S'$.\footnote{La page
écrit « Donc (réf.\ \ldots\ !!!) », une référence laissée à trouver. Le
résultat qu'elle demande est le lemme des pages 34 à 36 (l'algébricité d'un
schéma formel propre et lisse se vérifie sur la base réduite) ; ce
rapprochement est le nôtre, et ce lemme est énoncé pour une base locale
complète, non pour le complété le long d'un fermé quelconque.} Or
\[
  U \sqcup S' \longrightarrow S
\]
est fidèlement plat, et $A_U \sqcup A'$ est muni d'une donnée de descente
relativement à ce morphisme. Elle est effective : la question ne dépend que
de la réduction $S_0$ de $S$, où elle l'est, puisque $A_0$ existe --- c'est ce
que la page attribue à « la réduction de Raynaud ».\footnote{« Elle est
effective, puisqu'elle dépend de $S_0 \to S$, donc elle l'était déjà avant,
d'après la réduction de Raynaud. Ce qu'on gagne ! », lecture en partie
incertaine. La descente fidèlement plate des schémas abéliens n'est pas
toujours effective ; les résultats de M.~Raynaud sur les faisceaux amples et
la descente des schémas abéliens sont publiés en 1970 (\emph{Faisceaux amples
sur les schémas en groupes}, Lecture Notes 119), et nous ne savons pas quel
énoncé la page a en vue. Un carré $S_0 \leftarrow S'_0$ au-dessus de
$S \leftarrow S'$ est dessiné en marge.}

Un alinéa b), « réduction au cas $S$ local, $T$ réduit à un point fermé »,
est commencé puis biffé ; il se poursuit, biffé, en tête de la page 8.

\pagerange{8}{10}
\subsection*{b) Cas $S$ géométriquement unibranche, par l'espace de modules (pages 8 et 10)}

Après un revêtement convenable, $A_U$ admet une rigidification de Jacobi
--- une structure de niveau --- et une polarisation ; il définit donc un
morphisme $U \to R$ vers le schéma de modules $R$ des schémas abéliens
polarisés munis de cette structure.\footnote{« $R$ schéma modulaire à
décalage », lecture incertaine. L'existence d'un schéma de modules fin pour
les schémas abéliens polarisés avec structure de niveau $\ell \geq 3$ est le
théorème de Mumford (\emph{Geometric Invariant Theory}, 1965).} Soit $S'$
l'adhérence du graphe de $U \to R$ dans $S \times R$. Par le critère valuatif,
$S' \to S$ est propre : sur un trait $V$ qui envoie son point générique dans
$U$, le groupe $p$-divisible de $A_U$ se prolonge (c'est $M$), donc $A_U$ a
bonne réduction sur $V$ --- c'est le cas d'un trait, que la page attribue à
Raynaud et que traitent les pages 47 à 52.\footnote{« Le critère valuatif et
le résultat de Raynaud dans le cas d'un trait montrent que $S' \to S$ est
propre », lecture en grande partie incertaine. Le rapprochement avec le
tapuscrit est le nôtre ; la page 52 parle du « th.\ de Raynaud » pour un
énoncé de cette forme.}

Il veut ensuite montrer que $S' \to S$ est un isomorphisme ; par le
théorème de Tate, cela revient à prouver que $A_U$ se prolonge, quand
$A_U[p^\infty]$ se prolonge en $M$, de façon unique à isomorphisme unique
près. Par le Main Theorem de Zariski, il suffit de voir que les fibres
$S'_s$ sont finies. Soit $S''$ le normalisé de $S'$, supposé fini sur $S'$.
Par le théorème de Tate, le groupe $p$-divisible de l'image inverse sur $S''$
du schéma abélien universel est $M_{S''}$. Sur la fibre $S''_s$,
géométriquement connexe, cette famille a donc un groupe $p$-divisible
constant ; il en conclut qu'elle est elle-même constante à un morphisme
radiciel près, que $(S''_s)_{\mathrm{red}} \to R$ se factorise par un point,
et que $S''_s$ est radiciel sur $k(s)$.\footnote{Toute cette conclusion, en
tête de la page 10, est barrée d'un trait oblique, et se lit mal. Le
mécanisme est vraisemblable --- par Serre--Tate, une famille de schémas
abéliens en caractéristique $p$ dont le groupe $p$-divisible est constant
a une application de modules de différentielle nulle, ce qui n'en fait pas
une constante mais un morphisme radiciel ---, mais la page ne le démontre
pas.} Le cas où $S''$ n'est pas fini sur $S'$ se traiterait « par la méthode
essentiellement de [\;] p.\ 75--76 ».

Ce que ce raisonnement peut donner, si on le complète, est que $S' \to S$ est
propre, quasi-fini, donc fini, birationnel et à fibres radicielles : un
homéomorphisme universel. Ce n'est un isomorphisme que si $S$ est
normal.\footnote{Sur une courbe à point de rebroussement, géométriquement
unibranche, le normalisé est fini, birationnel et radiciel sans être un
isomorphisme. La page écrit « Je veux : prouver que c'est un isom. ». C'est
la raison d'être des alinéas c) et d) qui suivent : on traite d'abord le cas
normal, puis on descend le long d'un morphisme fini.}

\pagerange{10}{12}
\subsection*{c) Cas $S$ normal (pages 10 et 12)}

Soit $S$ normal et intègre. Choisissons un nombre premier $\ell \neq p$,
$\ell \geq 3$,\footnote{La page écrit « $\ell \neq p, 3$ », qu'on lit comme
la condition usuelle d'une structure de niveau rigidifiante ; la lecture est
la nôtre.} et un revêtement fini étale surjectif $U' \to U$ sur lequel
$A_{U'}$ est muni d'une rigidification de niveau $\ell$. Soit $S'$ le
normalisé de $S$ dans le corps des fonctions de $U'$ ; c'est un schéma normal
intègre fini sur $S$, qui prolonge $U' \to U$.\footnote{En marge, cerclé :
« inutile au fond ! », relié à « Normalisons », lecture incertaine.} Par b),
appliqué sur $S'$ qui est normal, on trouve sur $S'$ un schéma abélien $A'$
prolongeant $A_U \times_U U'$, avec $A'[p^\infty] \simeq M_{S'}$.

Soit $S'' = S' \times_S S'$. Le morphisme $S' \to S$, fini et dominant sur un
schéma normal, est universellement ouvert ; donc $U''$ est dense dans $S''$,
et la pleine fidélité (station 2), appliquée sur $S''$, munit $A'$ d'une donnée
de descente relativement à $S' \to S$.\footnote{« $S'' \to S$ est ouvert car
$S' \to S$ est universellement ouvert » ; c'est le théorème de Chevalley
(EGA IV 14.4.4), que la page ne nomme pas.} Le morphisme fini surjectif
$S' \to S$ est de descente effective pour les revêtements finis étales (SGA 1,
exposé IX) ; la donnée de descente sur $A'[\ell]$ fournit donc un revêtement
fini étale en groupes $N$ de $S$, avec $N|_U \simeq A_U[\ell]$. Un revêtement
fini étale surjectif $S_1 \to S$ qui trivialise $N$ trivialise
$A_U[\ell]$ ; on peut donc refaire la construction avec $S'$ \emph{étale}
sur $S$, donc plat. Comme $A'$ est projectif sur $S'$, la donnée de descente
est alors effective, et $A'$ descend en un schéma abélien $A$ sur $S$ qui
répond à la question.\footnote{La page cite « SGA 1 VIII » ; la descente des
revêtements étales le long d'un morphisme fini surjectif est, à notre
connaissance, dans l'exposé IX.
L'effectivité pour $A'$ projectif le long d'un morphisme fini étale demande
une polarisation compatible à la donnée de descente, qu'on obtient par
norme ; la page dit seulement « $A'$ étant projectif sur $S'$ ». Un ajout
encadré propose de remplacer l'hypothèse par : « pourvu qu'il existe $S' \to
S$ fini surjectif, avec $S'$ géom.\ unib. ».}

\pagerange{12}{14}
\subsection*{d) Cas $S$ réduit (pages 12 et 14)}

Soit $S$ réduit, et $S' \to S$ son normalisé, fini par hypothèse. Comme $S$ est
géométriquement unibranche aux points de $T$, $S' \to S$ y est radiciel, et
l'image inverse $U''$ de $U$ est dense dans $S'' = S' \times_S S'$. Par c), on
trouve sur $S'$ un schéma abélien $A'$, et la pleine fidélité pour $(S'',
U'')$ le munit d'une donnée de descente relativement à $S' \to S$. Un
morphisme fini et épimorphique --- ce qu'est $S' \to S$, $S$ étant réduit ---
est un morphisme de descente pour les schémas abéliens et pour les groupes
$p$-divisibles (ce sont les corollaires 2 et 3 de la page 24) ; il reste à
montrer que la donnée de descente est effective.\footnote{Le début de la page
14, sept lignes, est encadré et biffé de huit traits : une réduction au cas
local, une factorisation, « l'effectivité des descentes sur les \ldots\ de
FGA », un « cf plus bas ». C'est la question que reprennent les pages 26 à
38, sans que la page 14 y renvoie.} La page s'interrompt ici sur ce point.

\pagerange{14}{18}
\subsection*{e) Cas général, par recollement (pages 14 à 18)}

On peut supposer $S$ réduit (par a)). Pour tout $s \in S$, le cas local donne
sur $S(s) = \mathrm{Spec}\,\mathcal{O}_{S,s}$ une solution $A(s)$ du problème
induit, avec $U(s) = U \times_S S(s)$, $T(s) = T \cap S(s)$, $M(s)$. Par
unicité, les $A(s)$ forment un système cohérent : pour $t$ générisation de
$s$, $A(t) \simeq A(s) \times_{S(s)} S(t)$, avec la transitivité. Il faut un
lemme qui fasse de ce « système local » un schéma abélien sur $S$.

\textbf{Lemme} (page 16). Sur $S$ localement noethérien, un système local de
schémas abéliens provient d'un schéma abélien.\footnote{En marge, cerclé :
« Douteux » ; à la ligne suivante, « [\;] unicité ? ». La démonstration qui
suit est un brouillon dont l'ordre de lecture est incertain et qui
s'interrompt au bas du feuillet. On ne prétend pas que le lemme soit vrai
sous cette forme : l'argument naturel --- étendre chaque $A(s)$ à un voisinage
de $s$, puis recoller --- demande que les isomorphismes de recollement soient
uniques, ce qui suppose que toute composante irréductible du voisinage passe
par le point, comme le dit la page 18.} Une fois $A$ obtenu, on a
$A \times_S U \simeq A_U$, et un isomorphisme $A[p^\infty] \simeq M$ qui
induit les isomorphismes donnés sur les $S(s)$ ; il est donné sur $U$, et la
pleine fidélité le prolonge.

Le feuillet 18, qu'il numérote « 2 bis », donne l'outil du recollement :

\textbf{Lemme} (page 18). Soient $S$ noethérien réduit, $t \in S$, et $B$ un
schéma abélien sur $S$. Il existe un voisinage ouvert $V$ de $t$ tel que, pour
tout voisinage ouvert $V' \subset V$ de $t$, tout schéma abélien $B'$ sur $V'$
et tout isomorphisme $\varphi_t : B' \times_S S(t) \simeq B \times_S S(t)$, il
existe un isomorphisme $\varphi : B' \simeq B|_{V'}$ qui prolonge
$\varphi_t$, quitte à restreindre $V'$ ; il est unique dès que toute
composante irréductible de $V'$ contient $t$.\footnote{La phrase de l'énoncé
est en partie illisible (« il y a un isom.\ $\varphi(s)$ \ldots\ voisinage
\ldots\ unique \ldots »), et l'existence n'a de sens qu'après restriction du
voisinage ; c'est l'énoncé standard de prolongement d'un isomorphisme donné
sur l'anneau local (EGA IV 8). La page commence la démonstration par un
$S' \to S$ fini avec $S'$ réduit et géométriquement unibranche, et s'arrête sur
« Alors ».}

\pagerange{20}{20}
\section*{La reprise de la page 20}

La page 20 reprend le lemme de la page 3 pour $\ell = p$, avec une meilleure
conclusion. Soient $A$, $B$ schémas abéliens sur $S$ et
$v : A[p^\infty] \to B[p^\infty]$ ; \emph{quand $v$ provient-il d'un
homomorphisme $A \to B$ ?}

\begin{enumerate}
  \item[a)] Il y a un plus grand sous-schéma $S_v$ de $S$ sur lequel $v$ est
  algébrique ; il est fermé, et sa formation commute à tout changement de base
  localement noethérien.
  \item[b)] Si $t \in S_v$ et $k(t)$ est de caractéristique $p$, $S_v$
  contient un voisinage ouvert de $t$.
\end{enumerate}

\textbf{Corollaire.} Soit $Z$ l'ensemble des points associés de
$\mathcal{O}_S$. Supposons que pour tout $z \in Z$, l'adhérence
$\overline{\{z\}}$ rencontre $S_v$ en un point de caractéristique $p$. Alors
$v$ est algébrique.

\emph{Démonstration.} Par b), $S_v$ contient un ouvert $W$ qui rencontre
chaque $\overline{\{z\}}$ ; un ouvert étant stable par générisation, $W$
contient tous les points associés, et l'idéal $I$ de $S_v$ est nul en
chacun d'eux. Or les points associés de $I \subset \mathcal{O}_S$ sont des
points associés de $\mathcal{O}_S$ ; donc $I = 0$ et $S_v = S$.\footnote{La
démonstration est la nôtre ; la page n'en donne pas. Elle définit $Z$ comme
l'ensemble des points associés « \emph{et} » des points maximaux, qui en
font partie, et une incise biffée laisse un « $\neq p$ » isolé. Une première
version, biffée, cherchait le lieu $T$ fibre à fibre. Le corollaire de la
page 3 en est le cas où les points associés sont dans $U$.}

\pagerange{21}{24}
\section*{La descente par un épimorphisme propre (pages 21 à 24)}

Au pied d'un feuillet du tapuscrit anglais (page 21), quelques lignes de sa
main esquissent le schéma de toute la suite : $A$ sur $S$ et $A'$ sur $S'$,
deux images inverses $A''_1$, $A''_2$ sur $S''$ dont les groupes
$p$-divisibles sont isomorphes, $S \leftarrow S' \leftleftarrows S''$, et,
pour un $S'$ fait de deux copies, $S'' = S' \sqcup S'$.\footnote{Ces lignes
sont tête-bêche par rapport au texte dactylographié, et traversées de trois
traits courbes ; on n'en reprend que la structure.} La question est : un
schéma abélien sur $S'$ dont les deux images inverses sur $S''$ sont
identifiées descend-il à $S$ ?

Dans cette section, $f : S' \to S$ est un morphisme propre de schémas
localement noethériens, $S'' = S' \times_S S'$, $p_1, p_2 : S'' \to S'$ les
projections et $h = f p_1 = f p_2$. On dit que $f$ est un
\emph{épimorphisme} si $\mathcal{O}_S \to f_*\mathcal{O}_{S'}$ est injectif ;
propre et épimorphique, $f$ est surjectif.\footnote{C'est le N.B.\ de la page
22 ; la surjectivité, que nous ajoutons, vient de ce que l'image fermée de $f$
contient le support de $\mathcal{O}_S$.}

\textbf{Lemme 1} (page 22). Soient $f : S' \to S$ un épimorphisme propre de
schémas localement noethériens et $X$ un $S$-schéma formellement non ramifié.
La suite
\[
  \mathrm{Hom}_S(S, X) \longrightarrow \mathrm{Hom}_S(S', X)
  \rightrightarrows \mathrm{Hom}_S(S'', X)
\]
est exacte.\footnote{« propre » et « loc.\ » sont ajoutés en interligne ; en
marge : « [\;] inutile si $X$ loc.\ de t.f.\ ? ».}

\emph{Démonstration.} Soit $g' : S' \to X$ avec $g' p_1 = g' p_2$. Comme $f$
est surjectif et fermé, c'est une application quotient, et $g'$, constante
sur les fibres, se factorise en une application continue $g_0 : S \to X$. Le
morphisme $g'^{-1}\mathcal{O}_X \to \mathcal{O}_{S'}$ définit
$g_0^{-1}\mathcal{O}_X \to f_*\mathcal{O}_{S'}$, et l'égalité
$g' p_1 = g' p_2$ dit qu'il prend ses valeurs dans
\[
  \mathcal{A} = \mathrm{Ker}\bigl(f_*\mathcal{O}_{S'} \rightrightarrows
  h_*\mathcal{O}_{S''}\bigr).
\]
Donc $g'$ se factorise par $S' \to S'_1 = \mathrm{Spec}\,\mathcal{A} \to X$.
L'algèbre $\mathcal{A}$ est cohérente (EGA III 3, $f$ étant propre), donc
$S'_1 \to S$ est fini ; il est radiciel, parce que les deux morphismes
$S'' \rightrightarrows S'_1$ coïncident par définition de $\mathcal{A}$ et que
$S'' \to S'_1 \times_S S'_1$ est surjectif, si bien que $S'_1 \times_S S'_1$
est ensemblistement réduit à sa diagonale ; et c'est un épimorphisme, puisque
$\mathcal{O}_S \subset \mathcal{A} \subset f_*\mathcal{O}_{S'}$.\footnote{La
page écrit « $S'_1 \to S'$ fini radiciel », où la transcription signale que
le second $S'$ est peut-être un $S$ ; c'est $S'_1 \to S$ qui est fini et
radiciel. L'argument de radicialité est le nôtre ; la page dit « EGA III 3
indique aussitôt ».} On est donc ramené au cas d'un épimorphisme fini
radiciel :

\textbf{Corollaire 1} (page 22). Soient $f : S' \to S$ un épimorphisme entier
et radiciel, $S$ localement noethérien, et $X$ un $S$-schéma non ramifié,
localement de type fini. Alors $\mathrm{Hom}_S(S, X) \to \mathrm{Hom}_S(S',
X)$ est bijectif.\footnote{La ligne est très reprise, et l'ordre exact des
hypothèses incertain (« fini » biffé, « entier » en interligne, « ou $X \to
S$ loc.\ de t.f. »). La démonstration de la page 24 est presque entièrement
biffée ou illisible ; on en retient la réduction au cas local et un croquis
$\mathcal{O}_{S,s} \to \mathcal{O}_{X,x} \dashrightarrow \mathcal{O}_{S',s'}$.
Voici un argument : localement, $X$ est un sous-schéma fermé d'un $S$-schéma
étale $E$ ; un morphisme radiciel entier surjectif est un homéomorphisme
universel, donc $\mathrm{Hom}_S(S, E) = \mathrm{Hom}_S(S', E)$ (invariance
topologique du site étale, EGA IV 18.1.2) ; et si $\sigma : S \to E$
envoie $S'$ dans $X$, l'image réciproque $\sigma^{-1}(X)$ est un
sous-schéma fermé de $S$ par lequel $S' \to S$ se factorise, donc égal à $S$,
puisque $f$ est un épimorphisme.}

\textbf{Corollaire 2} (page 24). Un épimorphisme propre $f$ de schémas
localement noethériens est un morphisme de descente pour les schémas abéliens
relatifs : le foncteur image inverse est pleinement fidèle.

En effet, pour $A$, $B$ schémas abéliens sur $S$, le foncteur
$\underline{\mathrm{Hom}}_{\mathrm{gr}}(A, B)$ est représentable et non
ramifié sur $S$, et le lemme 1 s'applique à $X = \underline{\mathrm{Hom}}_
{\mathrm{gr}}(A,B)$.\footnote{La page le tient pour connu ; la
représentabilité est celle de FGA pour des schémas abéliens projectifs, la
non-ramification est la rigidité des homomorphismes.}

\textbf{Corollaire 3} (page 24). Le même $f$ est un morphisme de descente
pour les groupes $p$-divisibles.

L'argument du lemme 1 ramène à l'analogue du corollaire 1 pour
$X = \underline{\mathrm{Hom}}_{\mathrm{gr}}(A, B)$, $A$, $B$ groupes
$p$-divisibles, et il faut prouver que $X$ est formellement non ramifié. La
page s'arrête sur cette exigence, sans la démontrer.\footnote{Ce qu'il faut
est l'unicité du relèvement d'un homomorphisme de groupes $p$-divisibles le
long d'une immersion nilpotente ; c'est vrai (rigidité des homomorphismes de
groupes $p$-divisibles), mais n'est pas sur la page.}

\pagerange{26}{30}
\section*{L'effectivité : vers une base locale complète (pages 26 à 30)}

On suppose maintenant que $f$ est un épimorphisme \emph{effectif}, et l'on
veut montrer qu'il est de descente effective : un schéma abélien $A'$ sur
$S'$ muni d'une donnée de descente relativement à $S' \to S$ provient d'un
schéma abélien sur $S$.\footnote{« effectif » est d'une lecture incertaine.
Épimorphisme effectif s'entend au sens où $\mathcal{O}_S$ est le noyau de
$f_*\mathcal{O}_{S'} \rightrightarrows h_*\mathcal{O}_{S''}$, ce qu'utilise
l'exactitude formelle de la page 28.}

\emph{a) Réduction au cas $S$ local}, $S = \mathrm{Spec}\,R$ : par passage
à la limite.

\emph{b) Réduction au cas $S$ local complet} : on descend le long de
$\mathrm{Spec}\,\hat R \to \mathrm{Spec}\,R$, fidèlement plat, par le critère
de descente fpqc des schémas abéliens « de Raynaud ». Il faut alors savoir que
les morphismes déduits par changement de base de $f$ au-dessus de
$\mathrm{Spec}\,\hat R \times_S \mathrm{Spec}\,\hat R$, qui n'est plus
localement noethérien, sont encore de descente, et le corollaire 2 ne
s'applique plus. On s'en tire en remarquant que ce dont l'argument du lemme 1
avait besoin, que l'algèbre noyau
$\mathrm{Ker}\bigl(f_{2*}\mathcal{O}_{S'_2} \rightrightarrows
h_{2*}\mathcal{O}_{S''_2}\bigr)$ soit radicielle sur $\mathcal{O}_{S_2}$,
reste vrai au-dessus de $S_2$.\footnote{Le passage est en grande
partie illisible ; on en donne la structure. En marge, un tableau à quatre
lignes, $S \to S' \rightrightarrows S'' \equiv S'''$ et ses changements de
base au-dessus des étages $S_1$, $S_2$, $S_3$ du complexe de Čech de
$\mathrm{Spec}\,\hat R \to \mathrm{Spec}\,R$, chaque ligne étiquetée
« 1-descente » ou, pour la dernière, « 0-descente ». Les étiquettes ne sont
pas expliquées ; on les lit comme le degré de descente requis à chaque étage
(pleine fidélité, puis fidélité), et cette lecture est la nôtre. Un alinéa
c), « cas $\dim S = 0$ », est encadré et biffé.}

\pagerange{28}{30}
\subsection*{d) Construction d'un schéma abélien formel (pages 28 et 30)}

Soient $R$ local complet, $s$ son point fermé, $k$ son corps résiduel,
$S_0 = \mathrm{Spec}\,k$, et $S'_0$, $S''_0$ les fibres en $s$. D'abord,
$A'_0$ descend en une variété abélienne $A_0$ sur $k$.\footnote{La page dit
que $S'_0 \to S_0$ a une section, d'où un morphisme de descente effective
universelle ; la lecture de « section » est incertaine, et une section
demande un point rationnel. Sans point rationnel, on descend par un point
fermé de $S'_0$, d'extension résiduelle finie, le long de laquelle la
descente des variétés abéliennes (projectives) est effective.}

On complète $S'$, $S''$, $S'''$, $A'$, $A''$, $A'''$ le long des fibres en $s$
: on obtient un schéma abélien formel $\hat A'$ sur $\hat S'$ muni d'une
donnée de descente relativement à $\hat S' \to \hat S$, et il s'agit de
descendre dans la catégorie des schémas abéliens formels. Soient
$\Lambda_0 = \mathrm{Spf}\,W(k)$ et $\mathcal{M}$ le schéma formel des
modules de $A_0$.\footnote{Il note $T$ la base de Witt et $M$ l'espace des
modules formels ; on les renomme, $T$ et $M$ étant pris. L'anneau $W(k)$ des
vecteurs de Witt est la bonne base si $k$ est parfait ; sinon on prend un
anneau de Cohen de $k$.} Un lemme facile --- pratiquement trivial si $S'$ est
fini sur $S$ --- dit que se donner un prolongement formel $\hat A'$ de
$A'_0 = A_0 \otimes_k S'_0$ sur $\hat S'$ équivaut à se donner un
$\Lambda_0$-morphisme $\varphi' : \hat S' \to \mathcal{M}$, et de même sur
$\hat S''$. La donnée de descente sur $\hat A'$ qui prolonge celle de
$A'_0$ équivaut à $\varphi' \hat p_1 = \varphi' \hat p_2$. Or la suite
\[
  \hat S \longleftarrow \hat S' \leftleftarrows \hat S''
\]
est exacte dans les schémas formels --- trivial si $S'$ est fini sur $S$, et en
général par le théorème de comparaison formel-algébrique appliqué à
l'épimorphisme effectif $f$. Il existe donc $\varphi : \hat S \to \mathcal{M}$
avec $\varphi \hat f = \varphi'$, c'est-à-dire un schéma abélien formel $\hat
A$ sur $\hat S$ qui descend $\hat A'$.\footnote{Un croquis en marge empile
$\mathcal{M} \leftarrow S_0 \leftarrow S'_0 \leftleftarrows S''_0$ au-dessus
de $\Lambda_0 \leftarrow \hat S \leftarrow \hat S' \leftleftarrows \hat S''
\Lleftarrow \hat S'''$, avec les $\hat A$ au-dessus des colonnes. La page 26
avait annoncé cette construction comme « facile ».} Par le théorème de
comparaison --- le théorème d'existence de Grothendieck, EGA III 5 ---, il
reste à prouver que $\hat A$ est algébrique : si $\hat A$ est le complété d'un
schéma abélien $A$, l'isomorphisme formel $\hat A_{\hat S'} \simeq \hat A'$
s'algébrise et identifie $A_{S'}$ à $A'$ avec sa donnée de descente.

\textbf{Théorème} (page 30). Soient $R$ un anneau local noethérien complet,
$S = \mathrm{Spec}\,R$, et $\hat A$ un schéma abélien formel sur
$\hat S = \mathrm{Spf}\,R$. Les conditions suivantes sont équivalentes :
\begin{enumerate}
  \item[(i)] $\hat A$ est algébrisable : il existe un schéma abélien $A$ sur
  $S$ dont $\hat A$ est le complété formel ;
  \item[(ii)] il existe un morphisme propre surjectif $S' \to S$ tel que
  l'image inverse de $\hat A$ sur $\hat S'$ soit algébrisable ;
  \item[(iii)] pour tout quotient intègre $B$ de $R$ de dimension $1$, de
  normalisé $B'$, l'image inverse de $\hat A$ sur $\mathrm{Spf}\,B'$ est
  algébrisable.\footnote{La page écrit « quotient intègre $B$ de $A$ » ;
  c'est de $R$ qu'il s'agit. $B'$ est un anneau de valuation discrète
  complet, fini sur $B$.}
\end{enumerate}

\emph{(i) $\Rightarrow$ (ii)} est trivial. \emph{(ii) $\Rightarrow$ (iii).}
Soient $\Theta = \mathrm{Spec}\,B'$, un trait complet, et
$\Theta' = S' \times_S \Theta$, propre et surjectif sur $\Theta$, qui rend
$\hat A_{\hat\Theta}$ algébrique. Il existe un trait $\Theta_1$ fini sur
$\Theta$ et un $\Theta$-morphisme $\Theta_1 \to \Theta'$ ; donc
$\hat A_{\hat\Theta_1}$ est algébrique, et définit un schéma abélien
$A_{\Theta_1}$ sur $\Theta_1$, muni --- par le théorème de comparaison --- d'une
donnée de descente relativement à $\Theta_1 \to \Theta$. Elle est effective,
parce que $\Theta_1 \to \Theta$ est fini et fidèlement plat et que
$A_{\Theta_1}$, schéma abélien sur un trait, est
projectif.\footnote{Il note $T$ ce que nous appelons $\Theta$, et $T'$ son
changement de base ; la page ne redit pas que $T = \mathrm{Spec}\,B'$.}

\pagerange{32}{34}
\section*{Le théorème d'algébrisation : (iii) entraîne (i) (pages 32 et 34)}

\emph{(iii) $\Rightarrow$ (ii$'$)}, où (ii$'$) est (ii) avec $S' \to S$
\emph{fini} surjectif. Soient $S_n = \mathrm{Spec}\,R/\mathfrak{m}^{n+1}$,
$A_n$ la réduction de $\hat A$ sur $S_n$, $N_n =
\underline{\mathrm{NS}}_{A_n/S_n}$ et $P_n \subset N_n$ l'ouvert des
polarisations. Les $P_n$ se recollent en un schéma formel $\hat P$ sur $\hat S$,
réunion d'une famille dénombrable de schémas finis et non ramifiés sur $S$ :
\[
  \hat P = \coprod_\alpha \widehat{P^\alpha}, \qquad P^\alpha \text{ fini et
  non ramifié sur } S, \qquad Z_\alpha = \mathrm{Im}(P^\alpha \to S).
\]
L'hypothèse (iii) donne, pour chaque $B$, une polarisation de
$\hat A_{\mathrm{Spf}\,B'}$, donc un $S$-morphisme $\mathrm{Spec}\,B' \to
P^\alpha$ pour un $\alpha$ convenable : le point générique $x$ de
$\mathrm{Spec}\,B$ est dans $Z_\alpha$. La réunion $\bigcup_\alpha Z_\alpha$
contient donc tous les $x \in S$ tels que $\dim \overline{\{x\}} = 1$ (et le
point fermé). Par la proposition de la page 36, un nombre fini de
$Z_\alpha$ recouvrent $S$ ; la somme $S'$ des $P^\alpha$ correspondants est
finie et surjective sur $S$, et $\hat A_{\hat S'}$, muni de la polarisation
tautologique, est algébrisable par le théorème d'existence.

\emph{(ii$'$) $\Rightarrow$ (i)}, pour le $S'$ qu'on vient de construire. Par
le lemme sur $S_{\mathrm{red}}$ ci-dessous, on peut supposer $S$ réduit.
Quitte à faire une extension étale finie de $S$, et à redescendre --- ce qui
est possible « par Raynaud » ---, on peut supposer triviales les extensions
résiduelles de $S' \to S$. Alors chaque composante connexe $S'_i$ de $S'$ est
locale, et $S'_i \to S$ est une immersion fermée.\footnote{Pour un morphisme
fini quelconque, c'est faux (le normalisé d'une courbe à rebroussement, à
extension résiduelle triviale, n'est pas une immersion fermée). C'est vrai
ici parce que $S'$ est une somme de $P^\alpha$ non ramifiés sur $S$ : un
morphisme fini non ramifié vers un schéma local, d'extension résiduelle
triviale, est une immersion fermée (Nakayama). Les étiquettes de la page,
« (ii$'$) $\Rightarrow$ (i$'$) (cas $R$ réduit) » biffé puis
« (ii) $\Rightarrow$ (i) » plus bas, ne distinguent pas.} Si $S$ est
irréductible, l'une de ces immersions fermées est un isomorphisme, puisque
leur somme est surjective, et l'on conclut. Sinon, on passe aux composantes
irréductibles :

\textbf{Lemme} (page 34). Soient $R$ local noethérien complet, $S =
\mathrm{Spec}\,R$ réduit, et $\hat X$ un schéma formel propre sur $\hat S$. Si
pour toute composante irréductible $S_i$ de $S$, munie de sa structure
réduite, $\hat X \times_{\hat S} \hat S_i$ est algébrique, alors $\hat X$
l'est.

Par récurrence sur le nombre de composantes : $S = S' \cup S''$ avec $S'$,
$S''$ fermés, $\hat X_{\hat S'}$ et $\hat X_{\hat S''}$ proviennent de $X'$ et
$X''$, qui sont isomorphes au-dessus de $S' \cap S''$ par le théorème de
comparaison ; on recolle $X'$ et $X''$ par cet isomorphisme.\footnote{« On voit
facilement que $X$ recollé marche. » Comme $S$ est réduit, il est la somme
amalgamée de $S'$ et $S''$ le long de $S' \cap S''$ ; le recollement de $X'$
et $X''$ le long de leur partie commune est un pincement au sens de Ferrand
(2003), qui existe par exemple quand $X'$, $X''$ sont projectifs avec des
fibrés amples compatibles. Pour l'application aux schémas abéliens polarisés
de (ii$'$), c'est le cas ; la page ne pose pas la question.}

\pagerange{34}{36}
\section*{Algébricité modulo le nilradical (pages 34 et 36)}

\textbf{Lemme} (page 34). Soient $R$ local noethérien complet,
$S = \mathrm{Spec}\,R$, et $\hat X$ un schéma formel propre et lisse sur
$\hat S$. Si $\hat X \times_{\hat S} \hat S_{\mathrm{red}}$ est algébrique,
$\hat X$ l'est.\footnote{La page biffe « lisse » dans l'énoncé ; la
démonstration de la page 36 s'en sert. On le rétablit : l'énoncé ainsi
démontré suffit pour les schémas abéliens. Pour un $\hat X$ seulement propre
et plat, l'argument demanderait la théorie des déformations par le complexe
cotangent (Illusie, 1971), que les feuillets n'invoquent pas.}

On se ramène à un idéal $J$ de carré nul, $S_0 = V(J)$, avec
$\hat X_{\hat S_0}$ algébrique, de la forme $\hat X_0$ pour un $X_0$ propre et
lisse sur $S_0$. On compare la théorie des prolongements infinitésimaux de
$X_0$ en un $X$ lisse sur $S$ et celle de $\hat X_0$ en un $\hat X$ lisse sur
$\hat S$ :
\begin{enumerate}
  \item[1°)] l'obstruction vit dans « le même » groupe,
  \[
    H^2(X_0, \mathcal{T}_{X_0/S_0} \otimes_{\mathcal{O}_{S_0}} J) \simeq
    H^2(\hat X_0, \mathcal{T}_{\hat X_0/\hat S_0} \otimes \hat J),
  \]
  et c'est la même obstruction ;
  \item[2°)] les prolongements forment, de part et d'autre, un torseur sous
  « le même » groupe
  \[
    H^1(X_0, \mathcal{T} \otimes J) \simeq H^1(\hat X_0, \hat{\mathcal{T}}
    \otimes \hat J).
  \]
\end{enumerate}
Les isomorphismes sont ceux du théorème de comparaison, $R$ étant complet.
Comme le prolongement formel $\hat X$ existe, l'obstruction est nulle, il
existe un prolongement algébrique, et la complétion met en bijection les
classes de prolongements algébriques et formels ; $\hat X$ est donc le complété
d'un $X$.

\pagerange{36}{38}
\section*{Un argument de Baire (pages 36 et 38)}

Il reste à prouver la proposition dont (iii) $\Rightarrow$ (ii$'$) avait
besoin.

\textbf{Proposition} (page 36). Soient $R$ un anneau local noethérien
\emph{complet}, $S = \mathrm{Spec}\,R$, et $(Z_i)_{i \in I}$ une famille
dénombrable de parties fermées de $S$ dont la réunion contient tout point $x$
tel que $\dim \overline{\{x\}} \leq 1$. Alors une sous-famille finie recouvre
$S$.\footnote{La page biffe « anneau local noeth.\ complet » et écrit au-dessus
« schéma local (noethérien) ». Sans complétude, l'énoncé est faux : l'anneau
local de $\mathbb{Q}[x,y]$ en l'origine n'a qu'un ensemble dénombrable
d'idéaux premiers ; les adhérences de ses points de dimension $\leq 1$
forment une famille dénombrable de fermés propres qui contient tous ces
points, et aucune sous-famille finie ne recouvre le point générique. La
démonstration de la page utilise la complétude (théorème de structure de
Cohen) ; on la rétablit.}

Si $\dim S \leq 1$, c'est trivial : la réunion contient les points génériques,
qui sont en nombre fini. En général, il suffit, composante irréductible par
composante, de trouver un $Z_i$ qui contienne le point générique ; c'est la
contraposée du

\textbf{Corollaire} (page 38). Si les $Z_i$ sont rares dans $S$ intègre, il
existe $x \in S$, $\dim \overline{\{x\}} = 1$, avec
$x \notin \bigcup Z_i$.

\emph{Démonstration.} Par le théorème de structure de Cohen, $R$ est fini sur
un anneau $V[[T_1, \ldots, T_n]]$, $V$ anneau de valuation discrète complet
d'idéal maximal $\mathfrak{m}$.\footnote{La fin de la page 36, sous une grande
accolade barrée, est illisible ; on lit seulement l'existence de $V$ et
l'anneau $V[[T_1, \ldots, T_n]]$. Si $R$ contient un corps, on prend
$V = k[[T_0]]$.} Les images des $Z_i$ dans $\Theta = \mathrm{Spec}\,V[[T_1,
\ldots, T_n]]$ sont des fermés rares, et il suffit de trouver un point $x$ de
$\Theta$ de dimension $1$ hors de leur réunion : un point de $S$ au-dessus de
$x$ est alors de dimension $1$ et hors des $Z_i$. Les $V$-homomorphismes
$\varphi : V[[T_1, \ldots, T_n]] \to V$ sont paramétrés par
$(\varphi(T_1), \ldots, \varphi(T_n)) \in \mathfrak{m}^n$, espace métrique
complet ; le noyau de chacun définit un point $x_\varphi$ de dimension $1$.
Pour un fermé rare $Z \subset V(f)$, $f \neq 0$, la condition
$x_\varphi \in V(f)$ s'écrit $f(\varphi(T_1), \ldots, \varphi(T_n)) = 0$, et
définit une partie fermée rare de $\mathfrak{m}^n$. Par le théorème de Baire,
une réunion dénombrable de fermés rares de $\mathfrak{m}^n$ n'est pas
$\mathfrak{m}^n$ tout entier, et un $\varphi$ hors de cette réunion
convient.\footnote{« Il est immédiat que c'est une partie fermée rare » :
une série formelle non nulle ne s'annule identiquement sur aucun ouvert de
$\mathfrak{m}^n$. Le nom de Baire n'est pas sur la page.}

\pagerange{40}{44}
\section*{Le cas d'un schéma propre et plat (pages 40 à 44)}

\textbf{Remarque} (page 40). Soient $S = \mathrm{Spec}\,R$ local noethérien
complet de corps résiduel $k$, et $\hat X$ un schéma formel propre et plat sur
$\hat S$ tel que $H^0(X_0, \mathcal{O}_{X_0}) = k$. On suppose les
$\underline{\mathrm{Pic}}^0_{X_n/S_n}$ plats (par exemple
$\underline{\mathrm{Pic}}^0_{X_0/k}$ lisse), et l'on forme
\[
  \underline{\mathrm{NS}}_{X_n/S_n} = \underline{\mathrm{Pic}}_{X_n/S_n} /
  \underline{\mathrm{Pic}}^0_{X_n/S_n}, \qquad
  P_n = \underline{\mathrm{NS}}^+_{X_n/S_n} \ \text{(polarisations
  relatives)}.
\]
Comme aux pages 32 et 34, les $P_n$ se recollent en
$\hat P = \coprod_\alpha \hat P_\alpha$, les $P_\alpha$ finis sur $S$ et
connexes, et en procédant comme plus haut on trouve équivalentes :
\begin{enumerate}
  \item[(i)] il existe $S' \to S$ propre surjectif avec $\hat X_{\hat S'}$
  algébrisable et projectif ;
  \item[(ii)] pour tout sous-schéma fermé $T \subset S$ de dimension $\leq 1$,
  de normalisé $T'$, $\hat X_{\hat T'}$ est algébrisable et projectif ;
  \item[(iii)] il existe $S' \to S$ fini surjectif avec $\hat X_{\hat S'}$
  algébrisable et projectif ;
  \item[(iv)] il existe $S' \to S$ fini surjectif tel que, pour toute
  composante irréductible $S'_i$ de $S'$ munie de la structure réduite,
  $\hat X_{\hat S'_i}$ soit algébrisable et projectif.
\end{enumerate}
De plus, sous (iv), $\hat X$ est algébrique.\footnote{La page n'en dit pas
plus que « procédant comme plus haut ». La dernière phrase est lue « De plus,
\ldots\ (iv), $\hat X_{\hat S}$ est algébrique », après une rature ; une
dernière ligne, biffée, reprend (iv). Un N.B.\ en marge, relié à (iii) et
(iv), parle de $\underline{\mathrm{Pic}}_{X_0/k}$ lisse et d'un $S' \to S$
« fini et plat », et une note au bas de la marge d'« effectivité » ; ils se
lisent trop mal pour être restitués. Ni les équivalences ni la dernière
conclusion ne sont démontrées sur la page ; on les donne comme énoncés.}

\emph{N.B.} (page 42). Si $X_0$ est un schéma abélien, ces conditions
équivalent à ce que $\hat X$ soit algébrisable, et l'on peut supprimer la
condition de projectivité dans (i) à (iv).\footnote{C'est le théorème de la
page 30 ; pour (ii), un schéma abélien sur un trait est projectif. La page
ajoute en marge un « ? », et propose « ou seulement pro-abélien ? »,
lecture incertaine.}

\emph{Remarques de descente} (pages 42 et 44). Soient $S' \to S$ un morphisme
de descente effective, $S$ localement noethérien, et $X'$ propre et plat sur
$S'$, muni d'une donnée de descente relativement à $S$ ; la donnée est-elle
effective ? a) On se ramène aux anneaux locaux $\mathcal{O}_{S,s}$ ; b) par
descente fidèlement plate, à leurs complétés ; c) si les fibres de $X'$ n'ont
pas d'automorphismes infinitésimaux, on obtient un schéma formel $\hat X$ sur
$\hat S$ qui descend $\hat X'$ ; d) si, de plus, $X'$ devient projectif après
un $S'' \to S'$ surjectif et si les $\underline{\mathrm{Pic}}^0_{X'_n/S'_n}$
sont plats, la donnée de descente devient effective après une extension finie
étale $T$ de $\mathrm{Spec}\,\hat{\mathcal{O}}_{S,s}$, composante
irréductible par composante irréductible de $T$. Une dernière remarque énonce
: pour $S' \to S$ propre et effectif, $X'$ propre et plat sur $S'$ muni d'un
faisceau inversible $L'$, et une donnée de descente sur $(X', L')$, si aucune
fibre $(X'_{s'}, L'_{s'})$ n'a d'automorphisme infinitésimal, la donnée est
effective --- « argument par modules formels\ldots ».\footnote{Ces deux pages
sont un premier jet très rapide, surchargé ; une grande partie de la prose
qui relie les énoncés est illisible, et la moitié inférieure de la page 42
n'est donnée par la transcription qu'en fragments. On n'en restitue que les
énoncés lisibles, sans les tenir pour démontrés. L'absence d'automorphismes
infinitésimaux est ce qui rend les modules formels pro-représentables, et
l'argument annoncé est celui de la page 28, où $\mathcal{M}$ est l'espace des
modules formels.}

\pagerange{47}{50}
\section*{Le cas d'un trait : « Modèles de Néron et groupes $p$-divisibles » (pages 47 à 50)}

Ces quatre feuillets sont un tapuscrit, non signé, non daté, qui parle à la
première personne (« je sais maintenant prouver le théorème de Tate ») et cite
« la note sur les modèles de Néron, prop.\ 2.5 » et « le séminaire de l'an
dernier ». Ils sont annotés à l'encre bleue et au crayon.\footnote{Ni l'auteur
du tapuscrit ni la main des annotations ne sont établis par la transcription,
et on ne les établit pas ici. On donne les notes au crayon en note, comme
remarques d'un lecteur.}

Dans cette section, $S$ est un trait --- le spectre d'un anneau de valuation
discrète ---, de point générique $\eta$ et de point fermé $s$. Un $S$-schéma
en groupes $G$ est dit \emph{néronien} s'il est lisse, séparé, de type fini, et
possède la propriété universelle de Néron : tout $\eta$-morphisme
$Y_\eta \to G_\eta$ issu d'un $S$-schéma lisse $Y$ se prolonge en un
$S$-morphisme $Y \to G$ ; c'est le modèle de Néron de $G_\eta$.\footnote{Le
tapuscrit dit « $S$-néronien » sans définir le terme ; la définition est la
nôtre.} $H = (H_n)_n$ est un $S$-groupe $p$-divisible, $H_n$ le noyau de
$p^n$.

\textbf{Théorème.} Soient $G$ un $S$-schéma en groupes néronien, dont la fibre
générique est extension d'une variété abélienne par un tore, et $H$ un
$S$-groupe $p$-divisible. Tout homomorphisme $u_\eta : H_\eta \to G_\eta$ se
prolonge de manière unique en un $S$-homomorphisme $u : H \to
G$.\footnote{L'hypothèse sur la fibre générique est entre crochets au crayon,
et la marge dit « [\,] inutile », lecture incertaine. C'est ce que démontre
le lemme manuscrit des pages 51 et 52 ci-dessous.}

\textbf{Lemme 1.} Supposons $S$ hensélien. Soient $G$ néronien et commutatif,
$M$ un $S$-schéma en groupes fini et plat, $M^0$ sa composante neutre, et
$u_\eta : M_\eta \to G_\eta$ un homomorphisme. Pour que $u_\eta$ se prolonge,
il faut et il suffit que $u_\eta|_{M^0_\eta}$ se prolonge en
$u^0 : M^0 \to G$.\footnote{« commutatif » est ajouté, comme le demande au
crayon la marge : « $G$ commutatif ? ». Sans lui, $v_\eta$ ci-dessous n'est pas
un homomorphisme. Le tapuscrit note $\widetilde M$ la « partie connexe ».}

La nécessité est claire. Pour la suffisance, considérons
\[
  v_\eta : M_\eta \times G_\eta \longrightarrow G_\eta, \qquad (m, g)
  \longmapsto u_\eta(m) + g ;
\]
il suffit de prolonger $v_\eta$ en $v : M \times_S G \to G$, et $v|_M$
prolongera $u_\eta$. Soient $P_\eta = \mathrm{Ker}\,v_\eta$ et $P$ son
adhérence schématique dans $M \times_S G$. Comme $v_\eta$ identifie
$(M_\eta \times G_\eta)/P_\eta$ à $G_\eta$, la propriété de Néron prolonge
cette identification dès que le quotient $(M \times_S G)/P$ est lisse ; c'est
le

\textbf{Lemme 2.} Le $S$-schéma en groupes $(M \times_S G)/P$ existe et est
lisse sur $S$.\footnote{« existe et » est ajouté au crayon. Sur un trait, le
quotient d'un schéma en groupes de type fini par un sous-groupe fermé plat
existe ; c'est un résultat d'Anantharaman (\emph{Mém.\ SMF} 33, 1973), que
le tapuscrit ne pouvait citer sous cette forme.}

Le morphisme $v^0 : M^0 \times_S G \to G$, $(m, g) \mapsto u^0(m) + g$,
prolonge $v_\eta|_{M^0_\eta \times G_\eta}$, et il est scindé par
$g \mapsto (0, g)$ ; d'où une suite exacte
\[
  0 \to \mathrm{Ker}\,v^0 \to M^0 \times_S G \to G \to 0,
\]
où $\mathrm{Ker}\,v^0 \simeq M^0$ est plat et
$(M^0 \times_S G)/\mathrm{Ker}\,v^0 \simeq G$ lisse. Par ailleurs
$(M \times_S G)/(M^0 \times_S G) = M/M^0$ est étale, de sorte que
$(M \times_S G)/\mathrm{Ker}\,v^0$, extension d'un groupe étale par un groupe
lisse, est lisse. Enfin $\mathrm{Ker}\,v^0$, plat, est l'adhérence de sa fibre
générique, contenue dans $P_\eta$ ; donc $\mathrm{Ker}\,v^0 \subset P$, et
$(M \times_S G)/P$ est lisse comme quotient d'un groupe
lisse.\footnote{L'isomorphisme $\mathrm{Ker}\,v^0 \simeq M^0$, par
$m \mapsto (m, -u^0(m))$, est notre justification de « il en résulte que
$\mathrm{Ker}(\widetilde v)$ est plat ». En tête du feuillet 48, une note au
crayon, « en fait il y a \ldots\ : $\widetilde M$ ! », lecture très
incertaine.}

\emph{Démonstration du théorème.} L'unicité est claire, $H$ étant plat et $G$
séparé. La propriété de Néron se conserve par hensélisation : on peut supposer
$S$ hensélien. Par le lemme 1 appliqué aux $H_n$, on peut remplacer $H$ par sa
partie connexe, et supposer $H$ connexe.

\emph{a) Cas où $G$ est semi-stable}, c'est-à-dire où les fibres de la
composante neutre $G^0$ sont extensions de variétés abéliennes par des
tores. Soient $G_\eta[p^\infty]$ le groupe $p$-divisible de $G_\eta$ et $G'$
le $S$-groupe $p$-divisible défini par $G^0$, sa partie finie. On sait que
$G''_\eta = G_\eta[p^\infty]/G'_\eta$ est non ramifié, et se prolonge donc en
un $S$-groupe $p$-divisible étale $G''$.\footnote{« On sait » est entouré au
crayon, avec en marge : « “donc” abusif ». La correction à l'encre remplace
« non ramifié » par « étale ». C'est l'analogue $p$-adique de la filtration de
SGA 7, exposé IX, pour la réduction semi-stable (la « partie fixe » et le
quotient non ramifié) ; le tapuscrit ne cite rien.} Par le théorème de Tate,
le composé $H_\eta \xrightarrow{u_\eta} G_\eta[p^\infty] \to G''_\eta$ se
prolonge en $H \to G''$, qui est nul, $H$ étant connexe et $G''$ étale. Donc
$u_\eta$ se factorise par $G'_\eta$ et, encore par Tate, se prolonge en
$H \to G' \subset G$.

\emph{b) Cas général.} Il existe une extension finie et plate de traits
$S' \to S$, de point générique $\eta'$, telle que le modèle de Néron
$\mathcal{G}'$ de $G_{\eta'}$ sur $S'$ soit semi-stable.\footnote{Au crayon :
« ? Cas où $G_\eta$ est une V.A.\ seulement ! OK ». Le théorème de réduction
semi-stable (SGA 7, exposé IX) est énoncé pour les variétés abéliennes ; pour
une extension par un tore, il faut en outre déployer le tore, ce qui se fait
par une extension finie.} Par a), $u_{\eta'}$ se prolonge en
$u' : H_{S'} \to \mathcal{G}'$. Or $G$ se reconstruit à partir de
$\mathcal{G}'$ : la restriction de Weil
$\mathcal{G} = \prod_{S'/S} \mathcal{G}'$ est néronienne, $G_\eta$ s'identifie
à un sous-groupe de $\mathcal{G}_\eta$, et si $G_1$ est l'adhérence
schématique de $G_\eta$ dans $\mathcal{G}$, $G$ est le « lissifié » de
$G_1$.\footnote{« cf.\ note sur les modèles de Néron prop.\ 2.5 », texte que
nous n'identifions pas. Le lissifié est ce qu'on appelle aujourd'hui la
lissification (\emph{smoothening}) de Néron ; voir Bosch, Lütkebohmert et
Raynaud, \emph{Néron Models} (1990), § 3.1 et § 7.6.} Par la propriété
universelle de $\prod_{S'/S}$, $u'$ fournit $H \to \mathcal{G}$, qui prolonge
$u_\eta$ et se factorise par $G_1$, $H$ étant plat : soit
$u_1 : H \to G_1$. Il reste à relever $u_1$ à $G$. Comme $H$ est connexe, on
peut supposer $p$ égal à la caractéristique résiduelle de $S$.

\textbf{Lemme 3.} Soit $f : G \to G_1$ le morphisme canonique, isomorphisme
sur les fibres génériques. Il existe $n \geq 0$ et $g : G_1 \to G$ tels que
$f g = p^n$ dans $G_1$.\footnote{« ceci a été démontré au cours du séminaire
de l'an dernier ». La précision « qui induit isom.\ sur fibres gén.\ » est
une note au crayon.}

Alors $g u_1 : H \to G$ s'annule sur $H_n$, puisqu'il en est ainsi sur la fibre
générique ($g_\eta u_{1,\eta} = f_\eta^{-1} p^n u_{1,\eta}$) et que $H_n$ est
plat ; il se factorise donc par $p^n : H \to H$, de noyau $H_n$ : il existe
$u : H \to G$ avec $u \circ p^n = g u_1$. Alors
$f u \circ p^n = p^n u_1 = u_1 \circ p^n$, et comme $p^n$ est un épimorphisme,
$f u = u_1$ : $u$ prolonge $u_\eta$.\footnote{Le tapuscrit dessine le
diagramme à colonne exacte $0 \to H_n \to H \xrightarrow{p^n} H \to 0$, la
flèche $u$ en pointillé, et un arc $p^n$ de $G_1$ à $G_1$ au crayon. La
vérification $f u = u_1$ est la nôtre.}

\emph{Remarques du tapuscrit.} La démonstration utilise deux fois le théorème
de Tate, qui en toute rigueur n'est démontré que lorsque $\eta$ est de
caractéristique $0$. La première fois, le groupe d'arrivée $G''$ est étale, et
le théorème est alors facile. L'auteur sait maintenant le démontrer quand la
partie connexe du groupe d'arrivée, sur la fibre générique, est de type
multiplicatif.\footnote{Le théorème de Tate --- le foncteur « fibre
générique » est pleinement fidèle sur les groupes $p$-divisibles d'un anneau
de valuation discrète de caractéristique $0$ --- est de Tate (colloque de
Driebergen, 1966). Son extension en caractéristique $p$, sur une base normale,
est due à de Jong (\emph{Invent.\ Math.}\ 134, 1998). Tout le dossier, quand la
fibre générique est de caractéristique $p$, en dépend.}

\textbf{Corollaire.} Soit $G_\eta$ un groupe sur $\eta$, extension d'une
variété abélienne par un tore, et $p$ la caractéristique résiduelle de $S$,
supposée $> 0$. Pour que $G_\eta$ ait bonne réduction sur $S$, il faut et il
suffit que son groupe $p$-divisible se prolonge en un $S$-groupe
$p$-divisible.\footnote{« de car.\ rés.\ $p > 0$ » est une note au crayon ;
une autre, reliée à « bonne réduction », demande : « explicitez ce que ça veut
dire ». Le tapuscrit ne donne pas de démonstration et s'arrête là. Pour une
variété abélienne, c'est le critère $p$-adique de bonne réduction, analogue du
critère de Néron--Ogg--Šafarevič (Serre--Tate, \emph{Ann.\ of Math.}\ 1968,
pour $\ell$ distinct de la caractéristique résiduelle) ; c'est le cas d'un
trait du théorème de la page 2.}

\pagerange{51}{52}
\section*{Un lemme qui supprime l'hypothèse (pages 51 et 52)}

\textbf{Lemme} (page 51). Soient $K$ un corps, $G$ un groupe algébrique sur
$K$, $M = (M(n))_n$ un groupe $p$-divisible sur $K$ et $\varphi : M \to G$ un
homomorphisme. On suppose que $K$ est de caractéristique $p$, ou bien que $K$
est parfait et que $G$ ne contient pas de sous-groupe isomorphe à
$\mathbf{G}_a$. Alors $\varphi$ se factorise par un sous-groupe $H$ de $G$,
extension d'une variété abélienne par un tore.\footnote{L'hypothèse est un ajout
encadré, à la place d'un « b) » biffé ; on la lit comme une alternative,
« car $K = p$, ou [$K$ parfait et $G$ sans $\mathbf{G}_a$] », ce que
confirment les deux cas de la démonstration. « $p \neq$ car $K$ » est biffé
dans l'énoncé.}

\emph{Démonstration.} Soit $Z_n$ le centralisateur de $\varphi(M(n))$. Les
$Z_n$ forment une suite décroissante de sous-groupes fermés, donc
stationnaire ; soit $Z$ sa valeur finale. Les $\varphi(M(n))$ sont
commutatifs, contenus les uns dans les autres, et centralisent $Z$ : $\varphi$
se factorise par le centre de $Z$, qui est commutatif.\footnote{La page écrit
« par $Z$, et même par $\mathrm{Centr}(Z)$ » ; c'est par
$Z \cap \mathrm{Centr}(Z)$, le centre de $Z$, qu'il faut entendre la
factorisation, puisque c'est lui qui est commutatif.} On peut donc supposer
$G$ commutatif. Le composé $M \to G \to G/G^0$ est nul, un groupe fini ne
recevant pas d'image non nulle d'un groupe $p$-divisible ; on peut supposer
$G$ connexe.

\emph{1er cas : $K$ de caractéristique $p$.} Pour tout $\nu$, $\varphi$ se
factorise par $p^\nu G$, car $M(n) = p^\nu M(n + \nu)$. Pour $\nu$ grand,
$p^\nu G$ est extension d'une variété abélienne par un tore.\footnote{« à
dégager en lemme, valable pour tout groupe algébrique commutatif sur un corps
de car.\ $p > 0$ ». C'est vrai : sur une clôture parfaite, $G/G_{\mathrm{red}}$ est
infinitésimal donc tué par une puissance de $p$, d'où
$p^{2\nu} G \subset p^\nu G_{\mathrm{red}} \subset p^\nu G$, et ces groupes
coïncident pour $\nu$ grand ; la décomposition de Chevalley de
$G_{\mathrm{red}}$, extension d'une variété abélienne par $T \times U$,
montre que $p^\nu$ tue $U$ ; être une extension d'une variété
abélienne par un tore descend au corps de base. L'argument est le nôtre.}

\emph{2e cas : $K$ parfait, sans $\mathbf{G}_a$ dans $G$.} Alors
$G_{\mathrm{red}}$ est lisse, et sa partie unipotente, déployée sur un corps
parfait, est nulle faute de $\mathbf{G}_a$ : $G_{\mathrm{red}}$ est extension
d'une variété abélienne par un tore. Si $p$ est distinct de la
caractéristique, les $M(n)$ sont étales, et $\varphi$ se factorise par
$G_{\mathrm{red}}$ ; si $p$ est la caractéristique, on est dans le 1er
cas.\footnote{La page dit seulement : « Alors on sait que $G$ est ext.\ d'une
VA par un tore ! », ce qui n'est vrai que de $G_{\mathrm{red}}$ ; « $p \neq$ car
$K$, les $M(n)$ sont alors étales » est biffé. Deux lignes encadrées et
biffées ouvrent la page 52.}

\textbf{Application} (page 52). Soient $S$ un trait de corps des fonctions $K$
et de corps résiduel $k$, $G$ un $S$-schéma en groupes néronien, $M$ un
$S$-groupe $p$-divisible et $\varphi_K : M_K \to G_K$. Alors $\varphi_K$ se
prolonge, de façon évidemment unique, en $\varphi : M \to G$.

Si $p$ est distinct de la caractéristique résiduelle, $M$ est étale, et la
propriété de Néron suffit. Supposons $p = \mathrm{car}\,k$. Alors $K$ est de
caractéristique $p$, ou de caractéristique $0$ et donc parfait ; et $G_K$ ne
contient pas de $\mathbf{G}_a$, puisqu'il admet un modèle de Néron. Le lemme
donne $H_K \subset G_K$, extension d'une variété abélienne par un tore, et
$\psi_K : M_K \to H_K$ par lequel $\varphi_K$ se factorise. $H_K$ admet un
modèle de Néron $H$, et l'inclusion $H_K \to G_K$ se prolonge en $i : H \to
G$.\footnote{La page écrit le but de $i$ comme un $K$ ; on lit $G$. Le modèle
de Néron de $H_K$ est de type fini parce que $G_K$, qui a un modèle de Néron de
type fini, ne contient pas de $\mathbf{G}_m$, et $H_K$ non plus ; c'est le
critère d'existence de Bosch, Lütkebohmert et Raynaud (\emph{Néron Models},
§ 10.2), que la page ne pouvait citer.} Le théorème du tapuscrit, appliqué à
$H$ --- la page dit « le th.\ de Raynaud pour $H$ » ---, prolonge $\psi_K$ en
$\psi : M \to H$, et $\varphi = i \circ \psi$. « On gagne ! »

L'hypothèse du tapuscrit sur la fibre générique est ainsi superflue, comme le
disait la marge au crayon de la page 47 : il suffit que $G$ soit néronien. Et
c'est sous cette forme que le cas d'un trait sert la page 8.

\end{document}
